
\documentclass{article}
\usepackage{colortbl}
\usepackage{makecell}
\usepackage{multirow}
\usepackage{supertabular}

\begin{document}

\newcounter{utterance}

\centering \large Interaction Transcript for game `bbh', experiment `logical\_deduction\_three\_objects', episode 47 with Qwen3.6{-}27B.
\vspace{24pt}

{ \footnotesize  \setcounter{utterance}{1}
\setlength{\tabcolsep}{0pt}
\begin{supertabular}{c@{$\;$}|p{.15\linewidth}@{}p{.15\linewidth}p{.15\linewidth}p{.15\linewidth}p{.15\linewidth}p{.15\linewidth}}
    \# & \multicolumn{2}{c}{Player} && \multicolumn{2}{c}{Game Master} \\
    \hline

    \theutterance \stepcounter{utterance}  
    & & & \multicolumn{4}{p{0.6\linewidth}}{
        \cellcolor[rgb]{0.9,0.9,0.9}{
            \makecell[{{p{\linewidth}}}]{
                \texttt{\tiny{[P1$\langle$GM]}}
                \texttt{A logical deduction task which requires deducing the order of a sequence of objects.} \\
\\ 
\texttt{Q: The following paragraphs each describe a set of three objects arranged in a fixed order. The statements are logically consistent within each paragraph. In a golf tournament, there were three golfers: Amy, Eli, and Eve. Eve finished above Amy. Eli finished below Amy.} \\
\texttt{Options:} \\
\texttt{(A) Amy finished last} \\
\texttt{(B) Eli finished last} \\
\texttt{(C) Eve finished last} \\
\texttt{A: (B)} \\
\\ 
\texttt{Q: The following paragraphs each describe a set of three objects arranged in a fixed order. The statements are logically consistent within each paragraph. On a shelf, there are three books: a white book, a green book, and an orange book. The green book is to the right of the white book. The orange book is the rightmost.} \\
\texttt{Options:} \\
\texttt{(A) The white book is the leftmost} \\
\texttt{(B) The green book is the leftmost} \\
\texttt{(C) The orange book is the leftmost} \\
\texttt{A: (A)} \\
\\ 
\texttt{Q: The following paragraphs each describe a set of three objects arranged in a fixed order. The statements are logically consistent within each paragraph. On a shelf, there are three books: a red book, a gray book, and a white book. The white book is to the left of the gray book. The red book is the second from the left.} \\
\texttt{Options:} \\
\texttt{(A) The red book is the leftmost} \\
\texttt{(B) The gray book is the leftmost} \\
\texttt{(C) The white book is the leftmost} \\
\texttt{A: (C)} \\
\\ 
\texttt{Q: The following paragraphs each describe a set of three objects arranged in a fixed order. The statements are logically consistent within each paragraph. In a golf tournament, there were three golfers: Ana, Mel, and Ada. Ada finished above Ana. Mel finished below Ana.} \\
\texttt{Options:} \\
\texttt{(A) Ana finished first} \\
\texttt{(B) Mel finished first} \\
\texttt{(C) Ada finished first} \\
\texttt{A:} \\
            }
        }
    }
     \\ \\

    \theutterance \stepcounter{utterance}  
    & \multicolumn{4}{p{0.6\linewidth}}{
        \cellcolor[rgb]{0.561,0.737,0.561}{
            \makecell[{{p{\linewidth}}}]{
                \texttt{\tiny{[P1$\rangle$GM]}}
                \texttt{To determine the order of the golfers, let's analyze the statements provided:} \\
\\ 
\texttt{1.  **"Ada finished above Ana."**} \\
\texttt{    This means Ada's position is higher (better) than Ana's. In a ranking context, "above" usually implies a better rank (e.g., 1st is above 2nd). So, Ada > Ana.} \\
\\ 
\texttt{2.  **"Mel finished below Ana."**} \\
\texttt{    This means Mel's position is lower (worse) than Ana's. So, Ana > Mel.} \\
\\ 
\texttt{Combining these two relationships:} \\
\texttt{Ada > Ana > Mel} \\
\\ 
\texttt{This establishes the full order from first to last:} \\
\texttt{1.  Ada (First)} \\
\texttt{2.  Ana (Second)} \\
\texttt{3.  Mel (Last)} \\
\\ 
\texttt{Now let's evaluate the options:} \\
\texttt{(A) Ana finished first {-} Incorrect, Ada finished first.} \\
\texttt{(B) Mel finished first {-} Incorrect, Mel finished last.} \\
\texttt{(C) Ada finished first {-} Correct.} \\
\\ 
\texttt{A: (C)} \\
            }
        }
    }
    & & \\ \\

    \theutterance \stepcounter{utterance}  
    & & & \multicolumn{2}{p{0.3\linewidth}}{
        \cellcolor[rgb]{0.9,0.9,0.9}{
            \makecell[{{p{\linewidth}}}]{
                \texttt{\tiny{[GM$|$GM]}}
                \texttt{Target: (C)} \\
            }
        }
    }
    & & \\ \\

    \theutterance \stepcounter{utterance}  
    & & & \multicolumn{2}{p{0.3\linewidth}}{
        \cellcolor[rgb]{0.9,0.9,0.9}{
            \makecell[{{p{\linewidth}}}]{
                \texttt{\tiny{[GM$|$GM]}}
                \texttt{game\_result = LOSE} \\
            }
        }
    }
    & & \\ \\

\end{supertabular}
}

\end{document}
